\nonstopmode{}
\documentclass[a4paper]{book}
\usepackage[times,inconsolata,hyper]{Rd}
\usepackage{makeidx}
\makeatletter\@ifl@t@r\fmtversion{2018/04/01}{}{\usepackage[utf8]{inputenc}}\makeatother
% \usepackage{graphicx} % @USE GRAPHICX@
\makeindex{}
\begin{document}
\chapter*{}
\begin{center}
{\textbf{\huge Package `wcvpmatch'}}
\par\bigskip{\large \today}
\end{center}
\ifthenelse{\boolean{Rd@use@hyper}}{\hypersetup{pdftitle = {wcvpmatch: Taxonomic Name Reconciliation Against the 'WCVP' Backbone}}}{}
\ifthenelse{\boolean{Rd@use@hyper}}{\hypersetup{pdfauthor = {Paul Efren Santos Andrade}}}{}
\begin{description}
\raggedright{}
\item[Title]\AsIs{Taxonomic Name Reconciliation Against the 'WCVP' Backbone}
\item[Version]\AsIs{0.0.3}
\item[Author]\AsIs{Paul Efren Santos Andrade [aut, cre, cph]}
\item[Maintainer]\AsIs{Paul Efren Santos Andrade }\email{paulefrens@gmail.com}\AsIs{}
\item[Description]\AsIs{Standardizes and reconciles scientific plant names against
a World Checklist of Vascular Plants ('WCVP')-style taxonomic
backbone. The package parses names into taxonomic components and
applies staged exact and fuzzy matching for binomial and trinomial
inputs, including infraspecific rank-aware checks. It also returns
accepted-name context and row-level matching flags to support
reproducible, auditable preprocessing for downstream biodiversity,
spatial, and trait analyses. A user-supplied backbone can be passed
through 'target_df'; when the optional companion package 'wcvpdata'
is installed, its default checklist can also be used.}
\item[License]\AsIs{MIT + file LICENSE}
\item[Encoding]\AsIs{UTF-8}
\item[Roxygen]\AsIs{list(markdown = TRUE)}
\item[Imports]\AsIs{assertthat, cli, dplyr, fozziejoin, lifecycle, magrittr,
memoise, purrr, stringdist, stringr, tibble, tidyr}
\item[Depends]\AsIs{R (>= 4.1.0)}
\item[LazyData]\AsIs{true}
\item[LazyDataCompression]\AsIs{xz}
\item[Suggests]\AsIs{knitr, pkgload, rlang (>= 1.0.0), rmarkdown, testthat (>=
3.0.0), wcvpdata (>= 0.7.0)}
\item[VignetteBuilder]\AsIs{knitr}
\item[Additional\_repositories]\AsIs{}\url{https://paulesantos.r-universe.dev}\AsIs{}
\item[Config/testthat/edition]\AsIs{3}
\item[Config/Needs/website]\AsIs{pkgdown}
\item[URL]\AsIs{}\url{https://paulesantos.github.io/wcvpmatch/}\AsIs{,
}\url{https://github.com/PaulESantos/wcvpmatch}\AsIs{}
\item[BugReports]\AsIs{}\url{https://github.com/PaulESantos/wcvpmatch/issues}\AsIs{}
\item[Config/roxygen2/version]\AsIs{8.1.0}
\item[NeedsCompilation]\AsIs{no}
\end{description}
\Rdcontents{Contents}
\HeaderA{wcvpmatch-package}{wcvpmatch: Taxonomic Name Reconciliation Against the 'WCVP' Backbone}{wcvpmatch.Rdash.package}
\aliasA{wcvpmatch}{wcvpmatch-package}{wcvpmatch}
\keyword{internal}{wcvpmatch-package}
%
\begin{Description}
Standardizes and reconciles scientific plant names against a World Checklist of Vascular Plants ('WCVP')-style taxonomic backbone. The package parses names into taxonomic components and applies staged exact and fuzzy matching for binomial and trinomial inputs, including infraspecific rank-aware checks. It also returns accepted-name context and row-level matching flags to support reproducible, auditable preprocessing for downstream biodiversity, spatial, and trait analyses. A user-supplied backbone can be passed through 'target\_df'; when the optional companion package 'wcvpdata' is installed, its default checklist can also be used.
\end{Description}
%
\begin{Author}
\strong{Maintainer}: Paul Efren Santos Andrade \email{paulefrens@gmail.com} [copyright holder]

Authors:
\begin{itemize}

\item{} Paul Efren Santos Andrade \email{paulefrens@gmail.com} [copyright holder]

\end{itemize}


\end{Author}
%
\begin{SeeAlso}
Useful links:
\begin{itemize}

\item{} \url{https://paulesantos.github.io/wcvpmatch/}
\item{} \url{https://github.com/PaulESantos/wcvpmatch}
\item{} Report bugs at \url{https://github.com/PaulESantos/wcvpmatch/issues}

\end{itemize}


\end{SeeAlso}
\HeaderA{build\_genus\_index}{Build a Genus Index for Fast Prefiltering}{build.Rul.genus.Rul.index}
\keyword{internal}{build\_genus\_index}
%
\begin{Description}
\strong{[Stable]}

Creates a compact genus-level index from the target backbone. The index stores
one row per genus and a list-column with candidate \code{plant\_name\_id} values
associated with each genus.

If \code{plant\_name\_id} is not present in \code{target\_df}, a surrogate integer ID is
created to keep the index usable with custom backbones.
\end{Description}
%
\begin{Usage}
\begin{verbatim}
build_genus_index(target_df = NULL)
\end{verbatim}
\end{Usage}
%
\begin{Arguments}
\begin{ldescription}
\item[\code{target\_df}] Optional custom target table. If \code{NULL}, the optional \code{wcvpdata} checklist is used when available; otherwise pass a backbone explicitly.
\end{ldescription}
\end{Arguments}
%
\begin{Value}
A tibble with columns:
\begin{description}

\item[genus] Genus name (character).
\item[plant\_name\_id] List-column of unique IDs per genus.
\item[n\_records] Number of IDs per genus.
\item[genus\_nchar] Number of characters in the genus name.

\end{description}

\end{Value}
%
\begin{Examples}
\begin{ExampleCode}

target <- data.frame(genus = "Opuntia", species = "ficus-indica", plant_name_id = 1)
wcvpmatch:::build_genus_index(target)

\end{ExampleCode}
\end{Examples}
\HeaderA{classify\_spnames}{Classify Scientific Plant Names into Taxonomic Components}{classify.Rul.spnames}
%
\begin{Description}
\strong{[Stable]}

Parse and classify scientific plant names into taxonomic components:
genus, specific epithet, infraspecific rank, infraspecific epithet, and author.

Output is aligned to a backbone convention:
\begin{itemize}

\item{} \code{Orig.Genus} in Title Case (first letter uppercase, rest lowercase).
\item{} \code{Orig.Species} and \code{Orig.Infraspecies} epithets in lowercase.
\item{} \code{Infra.Rank} in lowercase (\code{subsp.}, \code{var.}, \code{subvar.}, \code{f.}, \code{subf.}).
\item{} \code{Author} is recovered from the input and preserved in its original
casing/punctuation (no forced uppercasing).
\item{} \code{Orig.Name} is reconstructed as: genus + species + (rank + infra) + author.

\end{itemize}


Robustness rules:
\begin{itemize}

\item{} \code{cf.} / \code{aff.} are removed from parsing but preserved as flags (\code{has\_cf}, \code{has\_aff}).
\item{} Hybrid markers (\code{x}/\AsIs{\texttt{\bsl{}u00D7}}) as standalone tokens are removed with \code{had\_hybrid = TRUE}.
\item{} \code{sp.} / \code{spp.} triggers genus-only classification (\code{Rank = 1}, \code{Orig.Species = NA})
and sets \code{is\_sp}/\code{is\_spp}.
\item{} If an infraspecific rank is present but the infraspecific epithet is missing,
sets \code{rank\_missing\_infra = TRUE} and keeps \code{Infra.Rank} while \code{Orig.Infraspecies = NA}.
\item{} If rank appears "late" (after author-like tokens), parsing is best-effort and
\code{rank\_late = TRUE}.
\item{} If there is no explicit rank and a third token exists, the function can infer an
unranked infraspecific epithet when the third token looks epithet-like (all lowercase),
and does not look like the start of an author. In that case \code{implied\_infra = TRUE},
\code{Orig.Infraspecies} is filled, \code{Infra.Rank = NA}, and \code{Rank = 3}.

\end{itemize}

\end{Description}
%
\begin{Usage}
\begin{verbatim}
classify_spnames(splist)
\end{verbatim}
\end{Usage}
%
\begin{Arguments}
\begin{ldescription}
\item[\code{splist}] Character vector. Scientific plant names.
\end{ldescription}
\end{Arguments}
%
\begin{Value}
A tibble with one row per input name and standardized columns/flags:
\begin{description}

\item[sorter] Numeric index of original order.
\item[Input.Name] Original input string as provided by user.
\item[Orig.Name] Reconstructed standardized name aligned to backbone + original-cased author.
\item[Orig.Genus] Genus in Title Case.
\item[Orig.Species] Specific epithet in lowercase, or \code{NA} for genus-only (\code{sp./spp.}).
\item[Author] Recovered author string (original casing/punctuation) or \code{""}.
\item[Orig.Infraspecies] Infraspecific epithet in lowercase (ranked or implied), or \code{NA}.
\item[Infra.Rank] Infraspecific rank in lowercase (\code{subsp.}, \code{var.}, \code{subvar.}, \code{f.}, \code{subf.}), or \code{NA}.
\item[Rank] Numeric level: \code{1} genus-only, \code{2} genus+species, \code{3} includes infraspecific epithet.
\item[has\_cf,has\_aff,is\_sp,is\_spp,had\_hybrid,rank\_late,rank\_missing\_infra,had\_na\_author,implied\_infra] Logical flags.

\end{description}

\end{Value}
%
\begin{Examples}
\begin{ExampleCode}
library(wcvpmatch)
classify_spnames(c("Opuntia sp.", "Rosa canina subsp. coriifolia (Fr.) Leffler"))
classify_spnames(c("Cydonia japonica tricolor")) # implied unranked infra epithet
\end{ExampleCode}
\end{Examples}
\HeaderA{direct\_match\_infra\_rank\_within\_species}{Direct Match Infraspecific Rank within Species}{direct.Rul.match.Rul.infra.Rul.rank.Rul.within.Rul.species}
\keyword{internal}{direct\_match\_infra\_rank\_within\_species}
%
\begin{Description}
Direct Match Infraspecific Rank within Species
\end{Description}
%
\begin{Usage}
\begin{verbatim}
direct_match_infra_rank_within_species(df, target_df = NULL)
\end{verbatim}
\end{Usage}
\HeaderA{fia}{Cleaned Master Tree Species List from FIA}{fia}
\keyword{datasets}{fia}
%
\begin{Description}
A cleaned dataset containing tree species recorded by the
Forest Inventory and Analysis (FIA) program of the U.S. Forest Service.
This dataset is used in the examples and README of the \code{wcvpmatch}
package. The data was downloaded in November 2022 from the official
webpage of the Forest Inventory and Analysis National Program, available
at the following \Rhref{https://research.fs.usda.gov/products/dataandtools/fia-datamart}{link},
and was originally used during the development of the \code{treemendous}
package. For \code{wcvpmatch}, the variable names have been standardized
to \code{Orig.Genus} and \code{Orig.Species}.
\end{Description}
%
\begin{Usage}
\begin{verbatim}
fia
\end{verbatim}
\end{Usage}
%
\begin{Format}
A data frame with 2169 rows and 2 variables:
\begin{description}

\item[Orig.Genus] Genus name of the species binomial
\item[Orig.Species] Specific epithet of the species binomial

\end{description}

\end{Format}
\HeaderA{fuzzy\_match\_infraspecies\_within\_species}{Fuzzy Match Infraspecific Epithet within Species}{fuzzy.Rul.match.Rul.infraspecies.Rul.within.Rul.species}
\keyword{internal}{fuzzy\_match\_infraspecies\_within\_species}
%
\begin{Description}
Fuzzy Match Infraspecific Epithet within Species
\end{Description}
%
\begin{Usage}
\begin{verbatim}
fuzzy_match_infraspecies_within_species(
  df,
  target_df = NULL,
  max_dist = 1,
  method = "osa"
)
\end{verbatim}
\end{Usage}
\HeaderA{prefilter\_target\_by\_genus}{Prefilter Target Backbone by Input Genera (Exact + Fuzzy)}{prefilter.Rul.target.Rul.by.Rul.genus}
\keyword{internal}{prefilter\_target\_by\_genus}
%
\begin{Description}
\strong{[Stable]}

Reduces the target backbone to genera relevant for the current input names.
This is designed as a pre-step before \code{wcvp\_matching()} to reduce search space.

Strategy:
\begin{itemize}

\item{} Exact genus candidates are always included.
\item{} Optional fuzzy genus candidates are included when \code{include\_fuzzy = TRUE}.
\item{} Returned object preserves the standard target schema used by the package.

\end{itemize}

\end{Description}
%
\begin{Usage}
\begin{verbatim}
prefilter_target_by_genus(
  df,
  target_df = NULL,
  genus_index = NULL,
  include_fuzzy = TRUE,
  max_dist = 1,
  method = "osa"
)
\end{verbatim}
\end{Usage}
%
\begin{Arguments}
\begin{ldescription}
\item[\code{df}] Input tibble/data.frame with either \code{Genus}/\code{Species} or \code{Orig.Genus}/\code{Orig.Species}.

\item[\code{target\_df}] Optional custom target table. If \code{NULL}, the optional \code{wcvpdata} checklist is used when available; otherwise pass a backbone explicitly.

\item[\code{genus\_index}] Optional pre-built index from \code{build\_genus\_index()}. If \code{NULL}, it is built on the fly.

\item[\code{include\_fuzzy}] Logical. If \code{TRUE}, include fuzzy-matched genera.

\item[\code{max\_dist}] Maximum fuzzy distance for genus matching (used when \code{include\_fuzzy = TRUE}).

\item[\code{method}] String distance method passed to \code{fozziejoin}.
\end{ldescription}
\end{Arguments}
%
\begin{Value}
A prefiltered \code{target\_df} tibble compatible with \code{wcvp\_matching(target\_df = ...)}.
Attributes:
\begin{description}

\item[candidate\_genera] Character vector of selected genera.
\item[exact\_genera] Character vector of exact matched genera.
\item[fuzzy\_genera] Character vector of fuzzy matched genera.

\end{description}

\end{Value}
%
\begin{Examples}
\begin{ExampleCode}

df <- data.frame(Genus = "Opuntia", Species = "yanganucensis")
target <- data.frame(genus = "Opuntia", species = "yanganucensis", plant_name_id = 1)
wcvpmatch:::prefilter_target_by_genus(df, target_df = target)

\end{ExampleCode}
\end{Examples}
\HeaderA{wcvp\_direct\_match}{Direct Match Species \& Genus Binomial or Trinomial names}{wcvp.Rul.direct.Rul.match}
\keyword{internal}{wcvp\_direct\_match}
%
\begin{Description}
\strong{[Stable]}

Tries to directly match Genus + Species | Genus + Species + Rank + Infraspecies to \AsIs{\texttt{WCVP data}}.
\end{Description}
%
\begin{Usage}
\begin{verbatim}
wcvp_direct_match(df, target_df = NULL)
\end{verbatim}
\end{Usage}
%
\begin{Arguments}
\begin{ldescription}
\item[\code{df}] \code{tibble} containing the species binomial split into the columns \code{Orig.Genus} and \code{Orig.Species}.

\item[\code{target\_df}] Optional custom target table. If \code{NULL}, the optional \code{wcvpdata} checklist is used when available; otherwise pass a backbone explicitly.
\end{ldescription}
\end{Arguments}
%
\begin{Value}
Returns a \code{tibble} with the additional logical column \code{direct\_match}, indicating whether the binomial was successfully matched (\code{TRUE}) or not (\code{FALSE}).
Returns original columns plus \code{Matched.Genus}, \code{Matched.Species}, \code{Matched.Infra.Rank}, and \code{Matched.Infraspecies}.
\end{Value}
%
\begin{Examples}
\begin{ExampleCode}

df_parsed <- classify_spnames("Opuntia yanganucensis")
target <- data.frame(genus = "Opuntia", species = "yanganucensis", plant_name_id = 1)
wcvpmatch:::wcvp_direct_match(df_parsed, target_df = target)

\end{ExampleCode}
\end{Examples}
\HeaderA{wcvp\_direct\_match\_species\_within\_genus}{Direct Match Species within Genus}{wcvp.Rul.direct.Rul.match.Rul.species.Rul.within.Rul.genus}
\keyword{internal}{wcvp\_direct\_match\_species\_within\_genus}
%
\begin{Description}
\strong{[Stable]}

Tries to directly match the specific epithet within an already matched genus in 'WCVP'.
\end{Description}
%
\begin{Usage}
\begin{verbatim}
wcvp_direct_match_species_within_genus(df, target_df = NULL)
\end{verbatim}
\end{Usage}
%
\begin{Arguments}
\begin{ldescription}
\item[\code{df}] \code{tibble} containing the species binomial split into the columns \code{Orig.Genus} and \code{Orig.Species}.

\item[\code{target\_df}] Optional custom target table. If \code{NULL}, the optional \code{wcvpdata} checklist is used when available; otherwise pass a backbone explicitly.
\end{ldescription}
\end{Arguments}
%
\begin{Value}
Returns a \code{tibble} with the additional logical column \code{direct\_match\_species\_within\_genus}, indicating whether the specific epithet was successfully matched within the matched genus (\code{TRUE}) or not (\code{FALSE}).
\end{Value}
\HeaderA{wcvp\_distribution}{Retrieve Tabular WCVP Distribution by Taxonomic Rank}{wcvp.Rul.distribution}
%
\begin{Description}
\strong{[Stable]}
\end{Description}
%
\begin{Usage}
\begin{verbatim}
wcvp_distribution(
  taxon,
  taxon_rank = c("species", "genus", "family", "order", "higher"),
  native = TRUE,
  introduced = TRUE,
  extinct = TRUE,
  location_doubtful = TRUE,
  wcvp_names = NULL,
  wcvp_distributions = NULL,
  prefilter_genus = TRUE,
  fallback_to_genus = TRUE,
  summarise_by_input = FALSE,
  max_dist = NULL,
  method = "osa",
  output = c("standard", "full", "spatial", "summary")
)
\end{verbatim}
\end{Usage}
%
\begin{Arguments}
\begin{ldescription}
\item[\code{taxon}] Character vector of taxa to query.

\item[\code{taxon\_rank}] Character scalar. One of \code{"species"}, \code{"genus"},
\code{"family"}, \code{"order"}, or \code{"higher"}. The last two require corresponding
\code{order} or \code{higher} columns in \code{wcvp\_names}.

\item[\code{native}] Logical. Include native occurrences? Defaults to \code{TRUE}.

\item[\code{introduced}] Logical. Include introduced occurrences? Defaults to
\code{TRUE}.

\item[\code{extinct}] Logical. Include extinct occurrences? Defaults to \code{TRUE}.

\item[\code{location\_doubtful}] Logical. Include doubtful occurrences? Defaults to
\code{TRUE}.

\item[\code{wcvp\_names}] Optional WCVP names table. If \code{NULL}, the function loads
\code{wcvpdata::wcvp\_matching\_names()}.

\item[\code{wcvp\_distributions}] Optional WCVP distribution table. If \code{NULL}, the
function loads \code{wcvpdata::wcvp\_distribution()}.

\item[\code{prefilter\_genus}] Logical. Forwarded to \code{wcvp\_matching()} for
species-level queries. Ignored for all other taxonomic ranks.

\item[\code{fallback\_to\_genus}] Logical. If \code{TRUE} and \code{taxon\_rank = "species"},
inputs without species-level distribution are retried at genus level.

\item[\code{summarise\_by\_input}] Logical. If \code{TRUE}, return one row per input taxon
with collapsed distribution fields. In this mode, \code{area\_codes}, \code{areas},
\code{regions}, \code{continents}, and \code{distribution} are returned as character
strings separated by \code{" - "} rather than list-columns.

\item[\code{max\_dist}] Maximum string distance. If \code{NULL}, species queries default
to \code{2} and genus/family queries to \code{0}.

\item[\code{method}] String distance method passed to \code{fozziejoin}.

\item[\code{output}] Output layout: \code{"standard"} (the default analytical
taxon-area table), \code{"full"} (the complete audit table), \code{"spatial"}
(the compact taxon-area table for a later spatial join), or \code{"summary"}
(one row per submitted taxon). \code{summarise\_by\_input = TRUE} is retained as
a backwards-compatible alias for \code{output = "summary"}.
\end{ldescription}
\end{Arguments}
%
\begin{Details}
Queries distribution records by matching a taxon name against the WCVP names
table and then resolving the corresponding rows in the WCVP distribution
table. The function is designed around \code{wcvpdata::wcvp\_matching\_names()} and
\code{wcvpdata::wcvp\_distribution()}, but custom tables with the same
schema can also be supplied.

Matching is performed with \code{fozziejoin}, using compact lookup tables and
length-based prefiltering to keep the candidate set small. Species queries are
resolved in two stages: genus candidates are matched first, then species
names are searched only within those candidate genera.

If species-level matches resolve to synonyms and the names table contains
\code{accepted\_plant\_name\_id}, distribution is recovered from the accepted taxon.
For queries above species, accepted names are preferred to avoid double
counting synonym records.

The result deliberately contains no geometry and does not require \code{sf}.
\code{area\_code\_l3} is the stable WGSrpd level-3 key intended for a later join
to a user-supplied spatial object.

Default names and distribution tables are cached for the current R session.
Exact species queries use a direct lookup; lowercase strings and species
keys over the full backbone are only built when a fuzzy query needs them.
\end{Details}
%
\begin{Value}
A non-spatial tibble. The default \code{output = "standard"} returns one
row per matched query-area combination with the submitted and resolved
taxa, geographic hierarchy, and four occurrence flags (12 columns).
\code{output = "full"}
additionally returns matching provenance and identifiers. \code{output = "spatial"} returns
only the taxon-area fields needed for a later spatial join. \code{output = "summary"} returns one row per input taxon with collapsed text fields such
as \code{distribution}, \code{areas}, \code{area\_codes}, \code{regions}, \code{continents}, and
\code{n\_areas}.
\end{Value}
%
\begin{Examples}
\begin{ExampleCode}


library(wcvpmatch)

wcvp_distribution("Opuntia ficus-indica", taxon_rank = "species")
wcvp_distribution("Opuntia", taxon_rank = "genus")
wcvp_distribution("Cactaceae", taxon_rank = "family")
# When `order` is present in a custom names table:
# \dontrun{wcvp_distribution("Caryophyllales", taxon_rank = "order", wcvp_names = custom_names)}


\end{ExampleCode}
\end{Examples}
\HeaderA{wcvp\_fuzzy\_match\_genus}{Fuzzy Match Genus Name}{wcvp.Rul.fuzzy.Rul.match.Rul.genus}
\keyword{internal}{wcvp\_fuzzy\_match\_genus}
%
\begin{Description}
\strong{[Stable]}

Tries to fuzzy match the genus name to the 'WCVP' table (using the optional \code{wcvpdata} checklist by default when available).
\end{Description}
%
\begin{Usage}
\begin{verbatim}
wcvp_fuzzy_match_genus(df, target_df = NULL, max_dist = 1, method = "osa")
\end{verbatim}
\end{Usage}
%
\begin{Arguments}
\begin{ldescription}
\item[\code{df}] \code{tibble} containing the species binomial split into the columns \code{Orig.Genus} and \code{Orig.Species}.

\item[\code{target\_df}] Optional custom target table. If \code{NULL}, the optional \code{wcvpdata} checklist is used when available; otherwise pass a backbone explicitly.

\item[\code{max\_dist}] Maximum edit distance used for fuzzy genus matching.

\item[\code{method}] String distance method passed to \code{fozziejoin} (for example \code{"osa"}).
\end{ldescription}
\end{Arguments}
%
\begin{Value}
Returns a \code{tibble} with the additional logical column \code{fuzzy\_match\_genus}, indicating whether the genus was successfully matched (\code{TRUE}) or not (\code{FALSE}).
Further, the additional column \code{fuzzy\_genus\_dist} returns the distance for every match.
\end{Value}
%
\begin{Examples}
\begin{ExampleCode}

df <- data.frame(Orig.Genus = "Opuntiaa", Orig.Species = "yanganucensis")
target <- data.frame(genus = "Opuntia", species = "yanganucensis", plant_name_id = 1)
wcvpmatch:::wcvp_fuzzy_match_genus(df, target_df = target)

\end{ExampleCode}
\end{Examples}
\HeaderA{wcvp\_fuzzy\_match\_species\_within\_genus}{Fuzzy Match Species within Genus}{wcvp.Rul.fuzzy.Rul.match.Rul.species.Rul.within.Rul.genus}
\keyword{internal}{wcvp\_fuzzy\_match\_species\_within\_genus}
%
\begin{Description}
\strong{[Stable]}

Tries to fuzzy match the species epithet within a matched genus against 'WCVP' (using the optional \code{wcvpdata} checklist by default when available).
\end{Description}
%
\begin{Usage}
\begin{verbatim}
wcvp_fuzzy_match_species_within_genus(
  df,
  target_df = NULL,
  max_dist = 1,
  method = "osa"
)
\end{verbatim}
\end{Usage}
%
\begin{Arguments}
\begin{ldescription}
\item[\code{df}] \code{tibble} containing the species binomial split into the columns \code{Orig.Genus} and \code{Orig.Species}.

\item[\code{target\_df}] Optional custom target table. If \code{NULL}, the optional \code{wcvpdata} checklist is used when available; otherwise pass a backbone explicitly.

\item[\code{max\_dist}] Maximum edit distance used for fuzzy species matching within genus.

\item[\code{method}] String distance method passed to \code{fozziejoin} (for example \code{"osa"}).
\end{ldescription}
\end{Arguments}
%
\begin{Value}
Returns a \code{tibble} with the additional logical column \code{fuzzy\_match\_species\_within\_genus}, indicating whether the specific epithet was successfully fuzzy matched within the matched genus (\code{TRUE}) or not (\code{FALSE}).
\end{Value}
\HeaderA{wcvp\_genus\_match}{Match Genus name}{wcvp.Rul.genus.Rul.match}
\keyword{internal}{wcvp\_genus\_match}
%
\begin{Description}
\strong{[Stable]}

Tries to match the genus name to the 'WCVP' table (using the optional \code{wcvpdata} checklist by default when available).
\end{Description}
%
\begin{Usage}
\begin{verbatim}
wcvp_genus_match(df, target_df = NULL)
\end{verbatim}
\end{Usage}
%
\begin{Arguments}
\begin{ldescription}
\item[\code{df}] \code{tibble} containing the species binomial split into the columns \code{Orig.Genus} and \code{Orig.Species}.

\item[\code{target\_df}] Optional custom target table. If \code{NULL}, the optional \code{wcvpdata} checklist is used when available; otherwise pass a backbone explicitly.
\end{ldescription}
\end{Arguments}
%
\begin{Value}
Returns a \code{tibble} with the additional logical column \code{genus\_match}, indicating whether the genus was successfully matched (\code{TRUE}) or not (\code{FALSE}).
\end{Value}
\HeaderA{wcvp\_matching}{Match Scientific Names Against WCVP}{wcvp.Rul.matching}
%
\begin{Description}
\strong{[Stable]}

Runs a matching pipeline with exact and partial matching for binomial and
trinomial names, including infraspecific rank validation.

When the default WCVP backbone is used, its normalized representation and
compact genus lookup are cached for the current R session. Prepared custom
backbones are also recognized inside the pipeline so internal matching nodes
do not normalize or deduplicate the same table repeatedly.
\end{Description}
%
\begin{Usage}
\begin{verbatim}
wcvp_matching(
  df,
  target_df = NULL,
  prefilter_genus = TRUE,
  allow_duplicates = FALSE,
  max_dist = 1,
  method = "osa",
  add_name_distance = FALSE,
  name_distance_method = "osa",
  profile = FALSE,
  output_name_style = c("snake_case", "legacy"),
  output = c("standard", "full")
)
\end{verbatim}
\end{Usage}
%
\begin{Arguments}
\begin{ldescription}
\item[\code{df}] Input tibble/data.frame with either \code{Genus}/\code{Species} or
\code{Orig.Genus}/\code{Orig.Species}. For trinomials, include \code{Infra.Rank} and
\code{Infraspecies} (or \code{Orig.Infra.Rank}/\code{Orig.Infraspecies}).

\item[\code{target\_df}] Optional custom target table. If \code{NULL}, data are read from
the optional \code{wcvpdata} checklist when available; otherwise pass \code{target\_df} explicitly.

\item[\code{prefilter\_genus}] Logical. If \code{TRUE}, prefilter \code{target\_df} to candidate
genera (exact + fuzzy) before running the matching pipeline.

\item[\code{allow\_duplicates}] Logical. If \code{TRUE}, duplicated taxon keys are
deduplicated internally for matching and then expanded back to original
rows. Output includes \code{input\_index} for traceability to the original input.

\item[\code{max\_dist}] Maximum distance used in all fuzzy matching stages (genus, species, infraspecies).

\item[\code{method}] A string indicating the fuzzy matching method (passed to \code{fozziejoin}). Supported methods:
\begin{itemize}

\item{} \code{"levenshtein"}: Levenshtein edit distance (default).
\item{} \code{"osa"}: Optimal string alignment.
\item{} \code{"damerau\_levensthein"} or \code{"dl"}: Damerau-Levenshtein distance.
\item{} \code{"hamming"}: Hamming distance (equal-length strings only).
\item{} \code{"lcs"}: Longest common subsequence.
\item{} \code{"qgram"}: Q-gram similarity (requires \code{q}).
\item{} \code{"cosine"}: Cosine similarity (requires \code{q}).
\item{} \code{"jaccard"}: Jaccard similarity (requires \code{q}).
\item{} \code{"jaro"}: Jaro similarity.
\item{} \code{"jaro\_winkler"} or \code{"jw"}: Jaro-Winkler similarity.
\item{} \code{"soundex"}: Soundex codes based on the National Archives standard.

\end{itemize}


\item[\code{add\_name\_distance}] Logical. If \code{TRUE}, add
\code{matched\_dist} as pairwise distance between input name
(\code{Input.Name} fallback \code{Orig.Name}) and \code{matched\_taxon\_name}.

\item[\code{name\_distance\_method}] Method passed to \code{stringdist::stringdist}
when \code{add\_name\_distance = TRUE} (for example \code{"osa"}).

\item[\code{profile}] Logical. If \code{TRUE}, attach a timing table in the
\code{"timings"} attribute of the returned tibble, with elapsed seconds per
pipeline stage.

\item[\code{output\_name\_style}] Naming style for output columns:
\begin{itemize}

\item{} \code{"snake\_case"} returns standardized lower snake\_case names.
\item{} \code{"legacy"} keeps the historical mixed naming convention.

\end{itemize}


\item[\code{output}] Output layout. \code{"standard"} (the default) returns the parsed,
matched, and accepted-name fields needed for routine reconciliation.
\code{"full"} additionally returns internal matching-stage flags, fuzzy
distances, and parsing diagnostics.
\end{ldescription}
\end{Arguments}
%
\begin{Value}
A tibble with parsed input fields and matched/accepted taxonomic
context. The standard output contains \code{input\_index}, input and matched
name components, authorship, matched and accepted IDs/names, status,
\code{match\_ambiguity}, and \code{matched}. Use \code{output = "full"} for matching
diagnostics.
\end{Value}
%
\begin{Examples}
\begin{ExampleCode}


library(wcvpmatch)
# Match a single name
wcvp_matching(data.frame(Genus = "Opuntia", Species = "yanganucensis"))

# Match multiple names with snake_case output
names <- c("Aniba heterotepala", "Anthurium quipuscoae")
df <- classify_spnames(names)
wcvp_matching(df, output_name_style = "snake_case")

# Attach per-stage timings for profiling
out <- wcvp_matching(df, output_name_style = "snake_case", profile = TRUE)
attr(out, "timings")


\end{ExampleCode}
\end{Examples}
\HeaderA{wcvp\_setup\_info}{Check Default Backbone Setup}{wcvp.Rul.setup.Rul.info}
%
\begin{Description}
\strong{[Stable]}
\end{Description}
%
\begin{Usage}
\begin{verbatim}
wcvp_setup_info(inform = TRUE)
\end{verbatim}
\end{Usage}
%
\begin{Arguments}
\begin{ldescription}
\item[\code{inform}] Logical. If \code{TRUE} (default), print a short setup message.
\end{ldescription}
\end{Arguments}
%
\begin{Details}
Reports whether the optional companion package \code{wcvpdata} is available for
use as the default WCVP backbone and, if not, explains how to install it
from \code{r-universe}.
\end{Details}
%
\begin{Value}
Invisibly returns a named list with setup status fields:
\code{default\_backbone\_available}, \code{wcvpdata\_installed},
\code{wcvpdata\_has\_backbone}, \code{wcvpdata\_version}, \code{repository}, and
\code{install\_command}.
\end{Value}
%
\begin{Examples}
\begin{ExampleCode}
library(wcvpmatch)
wcvp_setup_info()
\end{ExampleCode}
\end{Examples}
\HeaderA{wcvp\_suffix\_match\_species\_within\_genus}{Suffix Match Species within Genus}{wcvp.Rul.suffix.Rul.match.Rul.species.Rul.within.Rul.genus}
\keyword{internal}{wcvp\_suffix\_match\_species\_within\_genus}
%
\begin{Description}
\strong{[Stable]}

Tries to match the specific epithet by exchanging common suffixes within an already matched genus in 'WCVP'.
The following suffixes are captured: \code{c("a", "i", "is", "um", "us", "ae")}.
\end{Description}
%
\begin{Usage}
\begin{verbatim}
wcvp_suffix_match_species_within_genus(df, target_df = NULL)
\end{verbatim}
\end{Usage}
%
\begin{Arguments}
\begin{ldescription}
\item[\code{df}] \code{tibble} containing the species binomial split into the columns \code{Orig.Genus} and \code{Orig.Species}.

\item[\code{target\_df}] Optional custom target table. If \code{NULL}, the optional \code{wcvpdata} checklist is used when available; otherwise pass a backbone explicitly.
\end{ldescription}
\end{Arguments}
%
\begin{Value}
Returns a \code{tibble} with the additional logical column \code{suffix\_match\_species\_within\_genus}, indicating whether the specific epithet was successfully matched within the matched genus (\code{TRUE}) or not (\code{FALSE}).
\end{Value}
%
\begin{Examples}
\begin{ExampleCode}

df <- data.frame(Orig.Genus = "Opuntia", Orig.Species = "yanganucensa", Matched.Genus = "Opuntia")
target <- data.frame(genus = "Opuntia", species = "yanganucensis", plant_name_id = 1)
wcvpmatch:::wcvp_suffix_match_species_within_genus(df, target_df = target)

\end{ExampleCode}
\end{Examples}
\HeaderA{wcvp\_synonyms}{Retrieve Synonyms Resolved Through the WCVP Backbone}{wcvp.Rul.synonyms}
%
\begin{Description}
\strong{[Stable]}
\end{Description}
%
\begin{Usage}
\begin{verbatim}
wcvp_synonyms(
  taxon,
  target_df = NULL,
  prefilter_genus = TRUE,
  max_dist = 2,
  method = "osa",
  include_accepted = FALSE,
  output = c("compact", "full")
)
\end{verbatim}
\end{Usage}
%
\begin{Arguments}
\begin{ldescription}
\item[\code{taxon}] Character vector of species names to resolve.

\item[\code{target\_df}] Optional WCVP-like names table. If \code{NULL}, the optional
\code{wcvpdata} checklist is used.

\item[\code{prefilter\_genus}] Logical. Passed to \code{\LinkA{wcvp\_matching()}{wcvp.Rul.matching}}.

\item[\code{max\_dist}] Maximum string distance used while resolving \code{taxon}.

\item[\code{method}] String distance method passed to \code{\LinkA{wcvp\_matching()}{wcvp.Rul.matching}}.

\item[\code{include\_accepted}] Logical. Include the accepted name as a row with
\code{name\_role = "accepted"} before its synonyms? Defaults to \code{FALSE}.

\item[\code{output}] Output layout. \code{"compact"} (the default) returns the accepted
name and its synonym fields. \code{"full"} additionally returns matching
diagnostics, WCVP identifiers, basionym ID, and retrieval status.
\end{ldescription}
\end{Arguments}
%
\begin{Details}
Resolves submitted species names with \code{\LinkA{wcvp\_matching()}{wcvp.Rul.matching}} and retrieves every
WCVP record with \code{taxon\_status = "Synonym"} that points to the resolved
accepted name. This means that an accepted name, a synonym, or a minor
spelling variant can all be used as the query.

Exact submitted names use a direct lookup before invoking the fuzzy matching
pipeline. The accepted-name lookup and normalized default backbone are
cached for the current R session and shared with the other package functions.
\end{Details}
%
\begin{Value}
A tibble with one row per input name and retrieved synonym. The
default compact result contains \code{submitted\_name}, \code{accepted\_taxon\_name},
\code{accepted\_taxon\_authors}, \code{name\_role}, \code{synonym\_name},
\code{synonym\_authors}, and \code{homotypic\_synonym}. Inputs without a match or
without recorded synonyms are retained with missing synonym fields.
\end{Value}
%
\begin{Examples}
\begin{ExampleCode}


wcvp_synonyms("Nopalea cochenillifera")


\end{ExampleCode}
\end{Examples}
\printindex{}
\end{document}
